% formal/capacity-boundary-subdifferential.tex — subdifferential prediction (Q5 orbit-switching, Q5b exact boundaries)
%
% Extends [prop:capacity-smoothness-classification](b) from
% formal/capacity-smoothness-classification.tex
% to handle orbit appearance at feasibility boundaries.
%
% Sources:
%   - Historical empirical origin: old subdifferential-validation artifacts,
%     removed in the numerics replacement; use git history if needed.
%   - Supporting lemmas: formal/capacity-smoothness-classification.tex
%     (lem:orbit-feasibility-open, lem:per-orbit-smooth, lem:kkt-sensitivity)
%
% Structure:
%   §1 Directional derivative with orbit appearance (thm:subdiff-with-appearance)
%   §2 Open questions
%
% Labels defined here:
%   thm:subdiff-with-appearance


% ── §1 Directional derivative with orbit appearance ──
%
% Proposition [prop:capacity-smoothness-classification](b) assumes all tied
% orbits have beta > 0, so each A_sigma is C^1 on a full neighborhood of a_0.
% This fails when an orbit has beta_k = 0 for some k: the orbit is only
% feasible in the half-space of directions where beta_k does not go negative.
% The following theorem extends the directional derivative formula to this setting.

\begin{unverified}
\begin{theorem}[Directional derivative with orbit appearance]%
\label{thm:subdiff-with-appearance}
At~$a_0$, define the set of minimizing orbits
(feasible with minimum action):
\begin{equation}\label{eq:minimizer-set}
  \mathcal{R} = \bigl\{\sigma :
    A_\sigma(a_0) = c(a_0),\;
    \beta(\sigma; a_0) \geq 0,\;
    Q(\sigma; a_0) > 0,\;
    M_\sigma(a_0) \text{ non-singular}
  \bigr\}.
\end{equation}
Assume $|\mathcal{R}| \geq 1$ and that no orbit
outside~$\mathcal{R}$ achieves action~$c(a_0)$
with $\beta \geq 0$ and $Q > 0$.

For a perturbation direction~$d \in (\mathbb{R}^4)^F$,
an orbit $\sigma \in \mathcal{R}$ is
\emph{directionally feasible in direction~$d$} if for every
index~$k$ with $\beta_k(\sigma; a_0) = 0$,
\begin{equation}\label{eq:dir-feasible-condition}
  \nabla_a \beta_k(\sigma; a_0) \cdot d > 0.
\end{equation}
(The strict inequality ensures~$\beta_k$
becomes positive at first order under the perturbation.)
Here $\nabla_a \beta_k$ is the sensitivity from
Lemma~\ref{lem:kkt-sensitivity}.
The directionally feasible subset is
\begin{equation}\label{eq:dir-feasible-set}
  \mathcal{R}(d) = \bigl\{\sigma \in \mathcal{R} :
    \sigma \text{ is directionally feasible in direction } d
  \bigr\}.
\end{equation}
Orbits with $\beta(\sigma; a_0) > 0$ componentwise are
directionally feasible in every direction
(the condition~\eqref{eq:dir-feasible-condition} is vacuous).

Assume additionally that for every $\sigma \in \mathcal{R}$
and every~$k$ with $\beta_k(\sigma; a_0) = 0$,
$\nabla_a \beta_k(\sigma; a_0) \cdot d \neq 0$.
(This is a non-degeneracy condition; it holds for
generic~$d$ since $\nabla_a \beta_k \neq 0$ generically.)

If $\mathcal{R}(d) \neq \emptyset$, the one-sided directional
derivative of the capacity exists and equals
\begin{equation}\label{eq:subdiff-appearance}
  D_d^+\, c(a_0) = \min_{\sigma \in \mathcal{R}(d)}
    \nabla_a A_\sigma(a_0) \cdot d.
\end{equation}
\end{theorem}

\begin{proof}
Write $a(t) = a_0 + t\,d$ for $t \geq 0$.
We classify orbits in~$\mathcal{R}$ by their behavior
under the perturbation.

\emph{Step 1: Interior orbits remain feasible.}
If $\beta(\sigma; a_0) > 0$ componentwise, then
$\sigma \in \mathcal{R}(d)$ for all~$d$, and by
Lemma~\ref{lem:orbit-feasibility-open},
$\sigma$ remains feasible at~$a(t)$ for small~$t > 0$.
By Lemma~\ref{lem:per-orbit-smooth}, $A_\sigma$ is $C^\infty$
near~$a_0$, so
$A_\sigma(a(t)) = A_\sigma(a_0) + t\,\nabla_a A_\sigma(a_0) \cdot d + O(t^2)$.

\emph{Step 2: Directionally feasible boundary orbits remain feasible.}
Suppose $\sigma \in \mathcal{R}(d)$ and
$\beta_k(\sigma; a_0) = 0$ for some index~$k$.
By Lemma~\ref{lem:kkt-sensitivity}, the map
$a \mapsto \beta(\sigma; a) = M_\sigma(a)^{-1} b$
is smooth near~$a_0$ (the smoothness of the KKT solution holds
wherever $M_\sigma$ is non-singular, with no sign condition
on~$\beta$).
Therefore
\[
  \beta_k(\sigma; a(t))
  = \underbrace{\beta_k(\sigma; a_0)}_{=\,0}
  + t\,\underbrace{\nabla_a \beta_k(\sigma; a_0) \cdot d}_{>\, 0}
  + O(t^2)
  > 0 \quad \text{for small } t > 0,
\]
where the strict positivity of the first-order term uses
the directional feasibility condition~\eqref{eq:dir-feasible-condition}.
Components with $\beta_k(\sigma; a_0) > 0$ remain positive by
continuity.
Hence $\sigma$ is feasible at~$a(t)$ for small $t > 0$,
and the Taylor expansion of~$A_\sigma$ applies:
$A_\sigma(a(t)) = A_\sigma(a_0) + t\,\nabla_a A_\sigma(a_0) \cdot d + O(t^2)$.

\emph{Step 3: Directionally infeasible orbits become infeasible.}
If $\sigma \in \mathcal{R} \setminus \mathcal{R}(d)$,
there exists~$k$ with $\beta_k(\sigma; a_0) = 0$ and
$\nabla_a \beta_k(\sigma; a_0) \cdot d < 0$
(strict, by the non-degeneracy hypothesis).
Then $\beta_k(\sigma; a(t)) = t\,(\nabla_a \beta_k \cdot d) + O(t^2) < 0$
for small $t > 0$, so $\sigma$ is infeasible at~$a(t)$
and does not contribute to~$c(a(t))$.

\emph{Step 4: Upper bound.}
For every $\sigma \in \mathcal{R}(d)$,
$c(a(t)) \leq A_\sigma(a(t))
= c(a_0) + t\,g_\sigma \cdot d + O(t^2)$,
where $g_\sigma = \nabla_a A_\sigma(a_0)$.
Hence
\[
  c(a(t)) \leq c(a_0)
  + t\,\min_{\sigma \in \mathcal{R}(d)} g_\sigma \cdot d
  + O(t^2).
\]

\emph{Step 5: Lower bound.}
Any orbit~$\sigma$ contributing to $c(a(t))$ is feasible
at~$a(t)$ and thus either:
\begin{itemize}
\item
  belongs to $\mathcal{R}(d)$ (handled in Step~4), or
\item
  belongs to $\mathcal{R}$ but is directionally infeasible
  (impossible by Step~3: it is infeasible at~$a(t)$), or
\item
  lies outside~$\mathcal{R}$ with $A_\sigma(a_0) > c(a_0)$.
  By continuity of~$A_\sigma$,
  $A_\sigma(a(t)) > c(a_0) + \delta/2$ for some $\delta > 0$
  and small~$t$, so it cannot undercut the minimum.
% [TODO: JÖRN - this argument requires that every orbit sigma
%  outside R with beta(sigma; a_0) >= 0 has A_sigma(a_0) > c(a_0)
%  strictly. The third hypothesis in the theorem statement
%  ensures this. It also does not account for orbits that are
%  strictly infeasible at a_0 (some beta_k < 0, not just = 0)
%  but become feasible at a(t). Such orbits have A_sigma(a_0)
%  computed from the unconstrained KKT solution, which may be
%  below c(a_0). For the formula to hold, we need the action gap
%  between such orbits and c(a_0) to be bounded away from zero.
%  This is generically true but not formally proved here.]
\end{itemize}
Hence
\[
  c(a(t)) \geq c(a_0)
  + t\,\min_{\sigma \in \mathcal{R}(d)} g_\sigma \cdot d
  - O(t^2).
\]

\emph{Conclusion.}
Combining Steps~4 and~5:
\[
  c(a(t)) - c(a_0)
  = t\,\min_{\sigma \in \mathcal{R}(d)} g_\sigma \cdot d + O(t^2),
\]
so $D_d^+\, c(a_0)
= \lim_{t \to 0^+} \frac{c(a(t)) - c(a_0)}{t}
= \min_{\sigma \in \mathcal{R}(d)} g_\sigma \cdot d$.
\end{proof}
\end{unverified}


% ── §2 Open questions ──
%
% 1. Uniqueness of KKT solution for fixed sigma.
%    Lemma [lem:well-defined] shows Q = -xi is unique even when M is singular.
%    But when M is singular, the beta-components need not be unique, and the
%    LP search (rem:near-null-lp-search) selects one with beta > 0.
%    Question: is A_sigma(a) still smooth in a at points where M(a) is singular?
%    The action value is well-defined, but the IFT argument
%    (Lemma [lem:per-orbit-smooth]) requires non-singular M.
%
% 2. Lower bound gap in thm:subdiff-with-appearance.
%    The lower bound (Step 5) does not account for orbits that are strictly
%    infeasible at a_0 (beta_k(a_0) < 0 for some k) but become feasible
%    at a(t) with action below c(a_0). Such an orbit is NOT in R (since
%    R requires beta >= 0), and its unconstrained action A_sigma(a_0) may
%    be below c(a_0). For the formula to hold, we need A_sigma(a(t)) to
%    remain above c(a_0) + o(t) for such orbits. This requires the
%    infeasibility (negative beta_k) to be first-order, not merely
%    infinitesimal, so that A_sigma(a(t)) transitions smoothly and the
%    action gap is preserved. Generically true, but needs proof.
%
% 3. The nabla_a beta_k . d = 0 case.
%    The theorem avoids this by requiring strict inequality. For the
%    degenerate directions where nabla_a beta_k . d = 0, the formula
%    may still hold but requires second-order analysis. This is a
%    measure-zero set of directions (for generic a_0).
